\documentclass[journal]{IEEEtran}

\usepackage{amsmath,amssymb,amsfonts}
\usepackage{authblk}
\usepackage{graphicx}
\usepackage{booktabs}
\usepackage{array}
\usepackage{microtype}
\usepackage{xcolor}
\usepackage[ruled,linesnumbered]{algorithm2e}
\usepackage[hidelinks]{hyperref}

\graphicspath{{figures/}}

\newcommand{\revps}{\ensuremath{\,\mathrm{rev/s}}}
\newcommand{\fsenv}{\ensuremath{f_{\mathrm{env}}}}
\newcommand{\lockk}{\ensuremath{\mathrm{lock}_k}}
\newcommand{\snr}{\ensuremath{\mathrm{SNR}}}

% ---------------------------------------------------------------------------
% DRAFT MARKER.  \pending{...} flags a number or claim whose measurement has
% not yet landed.  Every such site states what will be measured.
% Remove this macro (and every use of it) before submission.
% ---------------------------------------------------------------------------
\newcommand{\pending}[1]{{\color{gray}\textsf{[pending:~#1]}}}

% Placeholder for a figure whose graphic is not yet generated.
\newcommand{\figplaceholder}[2]{%
  \framebox[\columnwidth]{\rule{0pt}{#1}\parbox{0.9\columnwidth}{\centering
      \pending{#2}}}}

\begin{document}

\title{Rotor Speed Estimation from Quadcopter Ego-Noise via Multi-Stage Iterative Phase-Noise Aware Vold-Kalman Filtering}

\author[1]{Dmitrii~Mukhutdinov \thanks{d.mukhutdinov@qmul.ac.uk}}
\author[1]{Lin~Wang \thanks{lin.wang@qmul.ac.uk}}

\affil[1]{School of Electronic Engineering and Computer Science, Queen
    Mary University of London}

\maketitle

\begin{abstract}
Noise emitted by unmanned aerial vehicles (drones) has a distinct harmonic
structure: for each rotor, both aerodynamic noise from the propeller and mechanical
noise from the motor are predominantly tonal, with fundamental frequencies propotional to
the rotor's variable rotational speed. An ability to reliably extract rotor speeds from
audio recording alone is useful: it can make drone tracking via microphone arrays more precise
and robust, it would enable motor behavior diagnostics when telemetry is not available and
post-annotation of audio-only drone noise data with motor speed trajectories -- the latter
enables more large-scale training of neural models for noise generation and suppression
conditioned on motor speeds data. Given how clearly structured the noise is, the problem
seems easy; however, in practice, when dealing with recordings
of a real quadcopter in free-flight conditions, it proves to be quite challenging.
The presence of four rotors with distinct, but close and intersecting rotational speed
trajectories, as well as the presence of significant broadband noise component and other
acoustic sources in the mixture make the problem too difficult to deal with for existing
methods for tacholess order tracking or multi-pitch estimation.
We show, however, that the problem is still feasible via an iterative trajectory search
and optimization method. Our proposed method include magnitude-only ridge tracking stage
for coarse initialization of speed trajectories and phase-noise aware iterative Kalman
filtering stage for their refinement. We demonstrate that the method is able to restore
the motor speeds from audio alone to near-telemetry precision, significantly outperforming
existing approaches. Moreover, the refinement stage is able to improve upon the imprecise
telemetry labels themselves. While the iterative algorithm is too computationally expensive to
be fit for real-time applications, we claim that it is very useful for pseudo-labeling
of ego-noise data.
\end{abstract}

\begin{IEEEkeywords}
Order tracking, Vold--Kalman filter, capture range, drone audition,
ego-noise, rotor speed estimation, instantaneous frequency, phase noise.
\end{IEEEkeywords}

\section{Introduction}
\label{sec:intro}

\IEEEPARstart{T}{his} version of the text is very very WIP. I really struggle with writing
by hand, apparently. This says a lot about my current state of mind and how well I understand
my own research. Funny. 2026...

What to say in introduction: we have very little free-flight drone noise data annotated with
motor speed telemetry -- basically it's just DREGON (we also have a few in-house
recordings -- Michael's data -- let's call it MD dataset). So it is really hard to train
quality neural models with such little data. But the task should be solvable - we mostly can
\emph{see} speed trajectories on spectrogram. But it is actually hard (reasons as in abstract).

Research questions:
\begin{enumerate}
\item Can we at least tell which trajectory fits the audio better?
\item Can we iteratively approach the true trajectories?
\end{enumerate}

Contributions:

\begin{enumerate}
\item We introduce harmonic-and-broadband phase model with independent phase noise components for every harmonic and justify it via phase coherence measurements on real data. 
\item A principled way of measuring fitness of the set of rotor speed trajectories to the audio, based on such noise model.
\item A coarse trajectories initialization stage based on Viterbi ridge-tracking (and maybe smth else)
\item A trajectories refinement stage which does Kalman filtering over instantaneous frequencies and achieves some real precision.
\end{enumerate}

\section{Related Work}
\label{sec:related}

\emph{VK order tracking and adaptive variants.} The VK filter was
introduced for high-slew order extraction~\cite{vold1993high}, extended to
multiple independent shafts~\cite{vold1997multi}, and characterized
in~\cite{herlufsen1999characteristics,tuma2000characteristics}; adaptive
reformulations followed~\cite{pan2007adaptive,pan2010extended}. Two recent
members split on the question this paper turns on: where the instantaneous
frequency (IF) is estimated. Li \emph{et al.}~\cite{li2020vkswt} estimate
IF \emph{outside} the filter, by synchrosqueezed-wavelet ridge detection,
and hand it to VK as a fixed parameter; a ridge detector has essentially
unbounded capture range, so this architecture never meets the failure we
analyze. IAVKF~\cite{jiang2024iavkf} estimates IF \emph{inside} the loop:
the envelope-phase derivative is smoothed by a second, IF-domain VK
problem and added to the carrier, while a per-component bandwidth adapts
to an orthogonality fixed point. IAVKF peels components greedily and
states it cannot handle overlapping IFs---precisely the quadrotor twin
pair---and its wide-to-narrow bandwidth adaptation carries no dependence
on harmonic order. Section~\ref{sec:vkloop} reads that adaptation as an
implicit capture schedule and makes it explicit.

\emph{Ridge extraction and tacholess order tracking.} Adaptive ridge
extraction~\cite{iatsenko2016ridges,aare2023ridge} and ridge-path
regrouping at crossings~\cite{chen2017rprg} are the blind front-end
competitors; tacholess angular-resampling methods recover the shaft rate
from the signal itself~\cite{peeters2019tacholess,wang2019tlot}, with
generalized demodulation~\cite{olhede2005generalized} and the fan-chirp
transform~\cite{weruaga2007fanchirp} as transform-domain relatives. These
methods target a single shaft; a single-component ridge detector cannot
prefer a shaft comb over its blade-passage octave or keep four
near-coincident combs apart, which is why our global stage scores whole
comb hypotheses rather than individual lines.

\emph{Drone audition.} Ego-noise suppression spans
BSS~\cite{wang2020bss}, learned masking~\cite{wang2020dnntf}, and deep
enhancement; several systems consume rotor
speeds~\cite{yen2023ncm,gulli2025rotorcond}. To our knowledge no prior
work recovers all four rates blind at twin separations below $1\revps$.

\section{Signal Model and Problem Statement}
\label{sec:model}

Audio $x_c(t)$ ($C \le 8$ microphones, $f_s = 16$\,kHz) is recorded on a
quadrotor whose rotor $i$ spins at $r_i(t)$ revolutions per second
(74--91\revps{} in cruise here). With shaft phase
$\varphi_i(t) = 2\pi \int_0^t r_i(\tau)\,d\tau$, the generative model is
\begin{equation}
  x_c(t) = \sum_{i=1}^{4} \sum_{k=1}^{K}
  \operatorname{Re}\!\left[ A_{c,i,k}(t)\, e^{\,\mathrm{j} k \varphi_i(t)}
  \right] + n_c(t),
  \label{eq:model}
\end{equation}
with slowly varying complex amplitudes $A_{c,i,k}$ and broadband noise
$n_c$. Real audio adds per-harmonic \emph{phase noise}: the phase of
harmonic $k$ is $\theta_k(t) = k\varphi_i(t) + b_k(t) + \psi_k$ with
$b_k$ a Brownian motion of diffusion rate $q_k$ (rad$^2$/s; coherence
time $\tau_k \sim 2/q_k$).

The primitive shared by all phase-based stages is \emph{demodulation}:
multiplying by a carrier built from the current track $\hat\varphi$ and
low-passing at bandwidth $B$ yields the complex envelope
\begin{equation}
  z_k(t) = \mathrm{LP}_B\!\big[ x(t)\, e^{-\mathrm{j} k \hat\varphi(t)}
  \big] \approx \tfrac{a_k}{2}\,
  e^{\,\mathrm{j}(k[\varphi-\hat\varphi] + b_k + \psi_k)} + n_k(t),
  \label{eq:envelope}
\end{equation}
sampled at an envelope rate $\fsenv \approx 62.5$\,Hz. A rate error
$\delta$ (\revps) rotates $z_k$ at $k\delta$\,Hz. For a signal-dominated
envelope the phase-noise law is
$\operatorname{Var}[\arg z_k] \approx 1/(2\,\snr_k)$, and because a rate
reading from harmonic $k$ divides the angular speed by $2\pi k$, the
Fisher information about $\delta$ scales as
$w_k \propto k^2 / \operatorname{Var}[\arg z_k] \propto k^2 |z_k|^2$---the
classical argument for favoring high harmonics, revisited empirically in
Sec.~\ref{sec:vkloop}.

Three per-stage quantities organize the design. \emph{Capture range}: the
largest initialization error from which a stage converges to the correct
mode. \emph{Resolution}: the accuracy it attains inside its basin.
\emph{Trajectory bandwidth}: the fastest rate-of-change it can represent.
A band of half-width $B$ retains harmonic $k$ only while $k|\delta| < B$,
so the recurring capture bound is $\delta_{\max} \approx B/k$ (in \revps),
halved when the criterion is the half-band.

\emph{Estimation target.} Recover the four trajectories
$\{r_i(t)\}_{i=1}^4$ from audio alone. Accuracy is scored as
permutation-invariant mean absolute error (PIT-MAE, \revps) against
telemetry: the track assignment that minimizes the pooled MAE.

\section{Method}
\label{sec:method}

Fig.~\ref{fig:pipeline} and Algorithm~\ref{alg:chain} summarize the
chain. The problem is conditionally linear: for fixed rates,
\eqref{eq:model} is linear in the amplitudes, and the amplitude block is
solved exactly (Sec.~\ref{sec:vksolve}). The rate block is multimodal at
every scale (octave doubles, twin swaps, ramps), so mode selection is done
by discrete global search in the magnitude domain (Sec.~\ref{sec:global})
and precision by a single local phase-domain pass
(Sec.~\ref{sec:pikalman}) whose capture range is matched to the error the
global stage leaves.

\begin{figure}[t]
  \centering
  \figplaceholder{7em}{F1: pipeline block diagram --- global comb DP
    $\to$ per-rotor decoupling $\to$ coupled VK envelope solve $\to$ one
    phase-increment ML pass}
  \caption{The blind chain. Each stage exists to hand its successor an
    error inside the successor's capture range. The VK \emph{frequency}
    loop is absent by design (Sec.~\ref{sec:vkloop}).}
  \label{fig:pipeline}
\end{figure}

\begin{algorithm}[t]
  \caption{Blind four-rotor rate tracking}
  \label{alg:chain}
  \KwIn{audio $x_c(t)$, $C$ channels, 16\,kHz}
  \KwOut{trajectories $\hat r_i(t)$, $i = 1..4$}
  \tcp{A. Global stage (magnitude domain)}
  seed bases: comb-matched scan of whitened mean spectrum, band cap
    1.2\,kHz, residual re-scan with span/dedup priors\;
  octave test: halve bases where line ratio $v(b)/v(2b) > 1.4$\;
  coarse full-range Viterbi (12--120\revps) on the rigid union template;
    trust gates; energy-timed takeoff bridge\;
  per-window ladder DP in $\pm6\revps$ around bases,
    \eqref{eq:combscore}--\eqref{eq:viterbi}\;
  \tcp{B. Per-rotor decoupling}
  \For{each rotor $i$ (coordinate descent)}{
    residual corridor DP, sibling teeth masked; accept only if surface
      quality $> 0.15$\;
  }
  residual re-seed of duplicated/faint rotors; truth-free per-rotor
    acceptance gate\;
  \tcp{C. Coupled VK envelope solve}
  solve \eqref{eq:normaleq} per coupling group (banded Cholesky,
    half-bandwidth $2M$); referee
    $\mathcal{R} = \|x - \hat x\|/\|x\|$\;
  \tcp{D. Phase-increment ML refinement (one pass)}
  \For{$B$-schedule $12 \to 2.5$\,Hz, $k$-caps coarse-to-fine}{
    demodulate \eqref{eq:envelope}; increments \eqref{eq:increment};
      estimate $q_k$ (moments); form $I[t], m[t]$ \eqref{eq:infoform};
      scalar Kalman $+$ RTS; twin gate; pair-joint two-peak fit on tight
      pairs\;
  }
  $\hat r_i \mathrel{+}= \delta r_i$ (clipped); \textbf{no} outer
    iteration\;
\end{algorithm}

\subsection{Global Stage: Comb-Matched Dynamic Programming}
\label{sec:global}

The first stage ignores phase. On a whitened log-magnitude spectrogram
$M(f,t)$ (running-median whitening over frequency, so combs stand out),
score a candidate shaft rate $c$ by the half-tooth contrast
\begin{equation}
  S(c,t) = \frac{1}{|\mathcal{K}|}\sum_{k \in \mathcal{K}}
  \Big( M(kc,\, t) - \big[M\big((k{-}\tfrac12)c,\, t\big)\big]_+ \Big),
  \label{eq:combscore}
\end{equation}
which penalizes sub-multiple aliases without the whitening-dip artifact of
a signed contrast, and chain frames by dynamic programming over a
discretized rate grid,
\begin{equation}
  V(c,t) = S(c,t) + \max_{c'} \big[ V(c', t{-}1) - \gamma\, |c - c'|
  \big].
  \label{eq:viterbi}
\end{equation}
The recursion returns the exact global maximizer over grid paths---no
initialization, no basins---so its capture range is the whole grid; its
resolution is capped by bin width and prior at $0.95$--$1.15\revps$ per
steady window.

\emph{Blind seeding.} Constant per-rotor bases come from a comb-matched
scan of the time-averaged whitened spectrum, capped at 1.2\,kHz (above
which aliases of lower fundamentals outscore true combs). Because a
strong rotor shadows a weaker one nearby, seeding proceeds in rounds:
each accepted comb is subtracted in the whitened domain and the residual
re-scanned. Acceptance is bounded by two plausibility priors---a span
bound (max/min ratio of the accepted set $\le 1.45$, with a sub-multiple
exemption) and mutual deduplication at $2.5\revps$---and the downstream
ladder is wrapped in a per-track guard that reverts any stage whose
output drops a track's comb \emph{occupancy} (fraction of harmonics above
the local noise floor); raw confidence rises during exactly the failure
the guard exists to catch.

\emph{Octave disambiguation.} Warmup spectra of two-bladed propellers
contain only even shaft harmonics, so comb evidence alone cannot decide
between the shaft comb and the blade-passage comb one octave up. The
discriminator is physical: if the line at base $b$ is stronger than the
line at $2b$ (median ratio over seed bases $> 1.4$), then $b$ is itself
the blade-passage comb and the bases are halved.

\emph{Ramps.} A coarse slope-tolerant Viterbi pass over the full physical
range (12--120\revps, $0.5\revps$ grid) on a short-window spectrogram
(2048-point FFT, 32\,ms hop) precedes the ladder, scoring the rigid union
template of all four bases moved together. A \emph{steady} gate (path
span $< 8\revps$) and a \emph{distrust} gate (path median departs the
seed median by $> 5\revps$) fall back to the constant init, and an
energy-timed bridge rebuilds the takeoff transition from the 50--200\,Hz
rumble-band energy profile when the narrowband evidence drowns in
spool-up noise.

\subsection{Per-Rotor Decoupling}
\label{sec:decouple}

The union DP estimates a common trajectory shape; individual rotors must
then be decoupled. A per-rotor \emph{residual corridor DP} runs
coordinate descent over rotors: for rotor $i$, the sibling combs' teeth
are masked out of the whitened surface and the DP re-solves in a corridor
($\pm8$ then $\pm4\revps$, grid $0.25$ then $0.1\revps$) around the
current track; the move is accepted only if the masked surface passes a
quality floor ($0.15$), since a corridor solve on a surface the siblings
dominate would re-capture their comb.

Duplicated or faint rotors---a blind scan can seed two rotors on one
physical comb and miss a fourth---are recovered by a residual
\emph{re-seed}: reconstruct the accepted combs, subtract, and re-scan the
residual for a replacement base. Every proposal passes a truth-free
per-rotor acceptance gate: decline when the sibling reconstruction
diverged, when the proposed mean move exceeds $0.5\revps$, or when the
proposal lands on a sibling's comb. On the real protocol the gate
declines 54 of 60 proposals, bounding the worst per-window loss at
$+0.019\revps$ against $+2.702$ ungated.

\subsection{Coupled VK Envelope Solve}
\label{sec:vksolve}

Fix the phase tracks $\hat\varphi_i(t)$. Model~\eqref{eq:model} is then
linear in the amplitudes, and the second-generation VK filter is
penalized least squares,
\begin{multline}
  \min_{\{A\}} \sum_t \Big( x(t) - \sum_{i,k}
  \operatorname{Re}\big[A_{i,k}(t)\, e^{\mathrm{j}k\hat\varphi_i(t)}\big]
  \Big)^{\!2} \\
  + \sum_{i,k} \rho_{i,k}^2 \sum_t \big| \Delta^2 A_{i,k}(t) \big|^2 ,
  \label{eq:vk2}
\end{multline}
where the second-difference penalty weight $\rho_{i,k}$ maps to a
$-3$\,dB analysis bandwidth by the standard
relation~\cite{tuma2000characteristics}. Indexing tracks by $m = (i,k)$
with carrier matrices
$C_m = \operatorname{diag}\big(e^{\mathrm{j}k\hat\varphi_i(t_n)}\big)$,
the normal equations are
\begin{equation}
  \big(C^{\mathsf H}C + R\big)\, a = C^{\mathsf H} x ,
  \label{eq:normaleq}
\end{equation}
with $R = \operatorname{blkdiag}(\rho_m^2 D_2^{\mathsf H} D_2)$. The
cross-track blocks $C_m^{\mathsf H}C_{m'}$ are diagonal, with unit-modulus
entries rotating at the instantaneous frequency gap between tracks;
\emph{time-major} interleaving of the unknowns places all couplings
within index distance $2M$, so the system is Hermitian positive-definite
banded with half-bandwidth $2M$, factorized by banded Cholesky in $O(N
M^2)$. Retaining the cross-track coupling is what lets two tracks
separated by less than a filter bandwidth divide energy correctly instead
of both claiming it---mandatory for interlaced quadrotor combs, and the
mechanism behind twin separation. Envelopes are narrowband by
construction, so the solve runs on carriers demodulated and decimated to
\fsenv; tracks are partitioned into coupling groups by union--find on
instantaneous separation, and far pairs are pruned.

This stage is the \emph{envelope engine} of the chain and the substrate
of the label-free referee $\mathcal{R}(r) = \|x - \hat x_r\|/\|x\|$
(Sec.~\ref{sec:setup}). Its classical companion, the VK
\emph{frequency}-update loop, is deliberately not used;
Sec.~\ref{sec:vkloop} shows why.

\subsection{Phase-Increment ML Refinement}
\label{sec:pikalman}

The one refinement pass measures the \emph{phase increment} between
consecutive envelope frames,
\begin{equation}
  \Delta\psi_k(t) = \arg\big( z_k(t)\, \overline{z_k(t{-}1)} \big),
  \label{eq:increment}
\end{equation}
wrapped, so a cycle slip corrupts at most one measurement. Under the
diffusion model, increments remain informative about the rate error
$\delta r$ long after absolute phase has decohered:
\begin{equation}
  \mathbb{E}[\Delta\psi_k(t)] = 2\pi\,\Delta t\, k\, \delta r(t),
  \label{eq:incmean}
\end{equation}
\begin{equation}
  \operatorname{Var}[\Delta\psi_k(t)] =
  \underbrace{\frac{1 - \operatorname{sinc}(2B\Delta t)}%
  {\snr_k}}_{\text{noise term}}
  + \underbrace{c_{\mathrm{diff}}\, q_k\, \Delta t}_{\text{diffusion
  term}} .
  \label{eq:incvar}
\end{equation}
\emph{Noise term.} Each phase reading has variance $1/(2\,\snr_k)$, but
the two readings share \emph{filtered} noise: after a low-pass of
bandwidth $B$ the noise autocorrelation at lag $\Delta t$ is
$\operatorname{sinc}(2B\Delta t)$, and the variance of a difference of
correlated variables gives the stated factor. At $B = 6$\,Hz and $\Delta
t = 16$\,ms, $1 - \operatorname{sinc}(0.192) \approx 0.06$: the increment
sees only 6\% of the per-sample noise variance, and omitting the factor
overpredicts noise ${\sim}16\times$ (and drives the $q_k$ estimate to
zero). \emph{Diffusion term.} The Brownian phase passes through the same
band before differencing; integrating its spectrum $q/(2\pi f)^2$ against
the differencing gain $2(1 - \cos 2\pi f\Delta t)$ over $|f| \le B$ gives
\begin{equation}
  \operatorname{Var} = q\,\Delta t \cdot \frac{2}{\pi}
  \int_0^{2\pi B \Delta t} \frac{1 - \cos u}{u^2}\, du
  = q\,\Delta t\; c_{\mathrm{diff}} ,
  \label{eq:cdiff}
\end{equation}
with $c_{\mathrm{diff}} \to 1$ as $B \to \infty$ (the unfiltered Brownian
increment) and $c_{\mathrm{diff}} \approx 0.19$ at our $B\Delta t$. The
diffusion rate $q_k$ is estimated from the data by method of moments: a
first pass with $q = 0$, then the robust excess of the residual increment
variance over the SNR-predicted noise term, divided by
$c_{\mathrm{diff}}\Delta t$.

Because every increment touches one time index, the per-frame evidence
collapses to two sufficient statistics,
\begin{equation}
  I[t] = \sum_{c,k} \frac{h_k^2}{R_{ck}[t]}, \qquad
  m[t] = \sum_{c,k} \frac{h_k\,\Delta\psi_{ck}[t]}{R_{ck}[t]},
  \label{eq:infoform}
\end{equation}
with gain $h_k = 2\pi k\,\Delta t$ and $R_{ck}$ from~\eqref{eq:incvar}.
Under a random-walk prior on $\delta r$ the MAP problem is a symmetric
tridiagonal system, solved by an information-form scalar Kalman forward
pass and Rauch--Tung--Striebel smoother---the forward recursion \emph{is}
forward elimination on that system and the smoother is the back
substitution. The in-band noise floor is estimated from an off-comb
demodulation ($+3$\,Hz off every comb line), and the information sum is
deflated by the effective-sample-size factor $2B\Delta t$, since the
62.5\,Hz frame rate oversamples the band.

\emph{Ambiguity gating.} An increment is unambiguous only while
$k\,|\delta r|\,\Delta t < \tfrac12$ cycle, so harmonics enter
coarse-to-fine ($k$-caps under a $B$-schedule $12 \to 2.5$\,Hz), and a
wrap guard drops near-$\pi$ increments. Rotors are refined sequentially
with a per-frame twin rule: measurements of harmonic $k$ of rotor $i$ are
gated out at frames where any harmonic of another rotor comes within $B +
B_{\mathrm{guard}}$ of it.

\emph{Twin pairs.} When two rotors sit within the band, the envelope is a
two-phasor sum, and its increments follow the \emph{stronger} tone
entirely: writing
\begin{equation}
  \arg z = \theta_1 + \operatorname{Im}
  \log\!\big(1 + \tfrac{a_2}{a_1}\, e^{\mathrm{j}(\theta_2 -
  \theta_1)}\big),
  \label{eq:winding}
\end{equation}
the log term has zero net winding for $a_2 < a_1$, so over a beat cycle
the average rotation rate of $z$ equals $\dot\theta_1$, biased by up to
half the split; no reweighting fixes it. Tight pairs are therefore
handled in \emph{pair-joint} mode: self-collided harmonics are
demodulated once by the pair-mean phase in a band wide enough for $\pm k
\cdot \mathrm{split}/2$; per short window a two-peak spectral fit
(parabolic interpolation, channel-incoherent power) yields two line
frequencies assigned to rotors \emph{by order} (higher line $\to$ faster
rotor), feeding the same per-rotor smoothers. The mixture's raw
increments are never used as a pair-mean measurement---that is the
content of~\eqref{eq:winding}.

The pass runs \emph{once}. Each application injects ${\approx}0.05\revps$
of fresh estimation noise that accumulates in the track, so iteration is
harmful (measured in Table~\ref{tab:iterharm}).

\subsection{Design Constants}
\label{sec:constants}

Table~\ref{tab:constants} freezes every constant of the chain; none is
tuned per recording.

\begin{table}[t]
  \centering
  \caption{Frozen design constants.}
  \label{tab:constants}
  \scriptsize
  \setlength{\tabcolsep}{4pt}
  \begin{tabular}{@{}ll@{}}
    \toprule
    Constant & Value \\
    \midrule
    Sample / envelope rate & 16\,kHz / 62.5\,Hz \\
    Seed scan band cap; dedup; span bound & 1.2\,kHz; $2.5\revps$; 1.45 \\
    Octave line-ratio threshold & 1.4 \\
    Coarse DP grid; step; penalty $\gamma$
      & 12--120\revps; $0.5\revps$; $0.4\revps$/frame \\
    Ladder DP corridor & $\pm6\revps$ around bases \\
    Corridor DP bands; grids; quality floor
      & $\pm8, \pm4\revps$; $0.25, 0.1\revps$; 0.15 \\
    Decoupling gate: max mean move & $0.5\revps$ \\
    $\pi$-Kalman $B$-schedule; harmonic range & $12 \to 2.5$\,Hz;
      $k \le 30$, coarse-to-fine \\
    Off-comb noise probe; update clip & $+3$\,Hz; $0.5\revps$ \\
    Refinement passes & 1 (no outer loop) \\
    \bottomrule
  \end{tabular}
\end{table}

\section{The VK Frequency Loop and Capture Range}
\label{sec:vkloop}

The classical VK refinement closes a loop over the envelope solve: read
each envelope's rotation rate, fuse over $k$ with the Fisher weights
$k^2|p_k|$, nudge $\hat r$, rebuild carriers, iterate. The envelope
solve, however, only reconstructs content within its bandwidth $B$ of the
assumed tooth position; a rate error $\delta$ displaces harmonic $k$ by
$k\delta$\,Hz, so once $k\delta \gtrsim B/2$ the true energy lies outside
the solved envelope and the loop's measurement is noise. The capture
radius is
\begin{equation}
  \delta_{\max} \approx \frac{B}{2k},
  \label{eq:basin}
\end{equation}
which at $B = 1.5$\,Hz and the informative harmonics $k \ge 6$ is below
$0.13\revps$---while the blind global stage delivers errors of
$0.55$--$1.5\revps$. At that operating point \emph{0 of 25} harmonics
have $k\delta$ in band anywhere in any window: the loop is inert, and,
being an in-band estimator, it does not fail loudly; it reports
convergence.

Three independent measurements confirm this. An instrumented per-factor
audit finds the fused slope RMS entering the smoother is
$0.003$--$0.008\revps$ against true errors of $0.55$--$1.5\revps$---a
shortfall of two to three orders of magnitude present before any gate or
shrinkage. An ablation with the loop disabled reproduces the full
pipeline to $\pm0.003\revps$ on every window. The label-free referee is
unchanged by five refine rounds to four decimal places ($0.5774 \to
0.5774$ on a real cruise window), while a phase-increment pass on the
same window improves it to $0.5300$. A phase-noise explanation is ruled
out: with the exact trajectory the envelope captures $100\%$ of comb
energy at up to $3\times$ the nominal shaft jitter, because VK
demodulates \emph{along the trajectory}, and correct demodulation removes
the frequency modulation by construction.

\emph{Relation to IAVKF.} IAVKF~\cite{jiang2024iavkf} contains this loop
mechanism for mechanism, but is initialized by wide-capture ridge
detection and adapts one scalar bandwidth per component, wide to narrow,
with \emph{no} dependence on harmonic order---an implicit capture
schedule, never analyzed as one (the paper's advice to widen the initial
band when the initial IF is poor is that schedule stated operationally).
With few low-order components, \eqref{eq:basin} does not bind; with 25
harmonics of four combs it collapses as $1/k$ on exactly the harmonics
the weighting favors.

Two principled fixes follow. First, the \emph{$k$-scaled bandwidth}
\begin{equation}
  B_k = k\, B_0
  \label{eq:kscaled}
\end{equation}
makes the capture radius $B_0$ at every harmonic. The separation argument
survives: two rotors' $k$-th teeth are $k|r_i - r_j|$ apart, so a
$k$-scaled filter preserves exactly the discrimination the fundamental
had. Coverage demands it too: at a fixed band the twin-collision gate
starves the comb---of 30 harmonics, DREGON keeps $20.5$ at $B = 1.5$\,Hz,
$15$ at $6$\,Hz, and \emph{0} at $12$\,Hz on every unit---while
$k$-scaling holds the collision radius constant. Second, the
\emph{measured weight law}: on 36 DREGON and 52 Matrice~100 measurement
units the empirically optimal weight is $1/v_k \propto k^{2.0}$ (indoor)
and $k^{1.5}$ (outdoor) with \emph{no} amplitude factor
(Fig.~\ref{fig:weightlaw}); since $|p_k| \propto k^{-2}$ on these
recordings, the deployed $k^2|p_k|$ is effectively \emph{flat} in $k$
where it should rise as $k^2$. Under $B_k = 0.25k$ the fitted exponent
collapses to ${\approx}0$ (equal information per harmonic); the same
measurement supports neither a variance floor nor a cap on $k \le 30$.
Whether these two fixes revive the frequency loop is an experimental
question, answered in Sec.~\ref{sec:revival} \pending{stage-C/D
iteration sweep, arms: control / $k$-scaled / IAVKF-adaptive / both /
both$+$measured weights, offsets 0/0.3/1.0\revps}.

\section{Experimental Setup}
\label{sec:setup}

\subsection{Synthetic Tiers}
Three synthesis tiers isolate error sources; each comprises 6 windows of
16\,s with frozen seeds. \textbf{S1}: one rotor, $K = 30$ harmonics with
the DREGON-measured amplitude rolloff, locked phases ($q_k = 0$),
broadband noise at $+10$\,dB, band-limited (8\,Hz) Ornstein--Uhlenbeck
shaft trajectory---isolates the estimator floor. \textbf{S2}: four rotors
including one twin pair at $0.5\revps$ separation, same locked-phase
model, 0\,dB---isolates inter-rotor interference. \textbf{S3}: four
rotors with per-harmonic phase diffusion at the measured $q_k$ law and
stronger broadband noise ($-5$\,dB)---decoherence robustness. A fourth
rung, \textbf{S4}, is not an estimation tier but the per-harmonic
variance measurement behind the weight law of Sec.~\ref{sec:vkloop}
(Fig.~\ref{fig:weightlaw}).

\subsection{Real Data}
\textbf{DREGON}~\cite{strauss2018dregon}: indoor free flight, 8-mic
array, resampled $44.1 \to 16$\,kHz, 1\,kHz motor telemetry; cruise twin
pairs at $0.43$ and $0.81\revps$. \textbf{Matrice~100} (recordings
FLY124/FLY125): outdoor, 8-mic ring, flight-controller telemetry at
${\approx}29$\,Hz. Evaluation uses a frozen 15-window protocol spanning
steady cruise, ramps, and warmup. The outdoor telemetry required
auditing: a label-free scan finds a linearly drifting clock ($R^2 =
0.94/0.92$; residual RMS $2.9/4.5$\,ms against $12/16$\,ms for the best
constant lag) and a $0.70\%/0.71\%$ rate scale, two independent fits
agreeing to $0.008$ percentage points. Corrected labels are used
throughout; the correction alone improves pooled cruise error by
$8.5\%$---label quality, not estimator credit.

\subsection{Controlled-Imprecision Initialization}
\label{sec:initproto}
Refiner capture ranges can only be compared from identical
initializations of controlled quality. We initialize
\begin{equation}
  r^{\mathrm{init}}_i(t) = r^{\mathrm{true}}_i(t) +
  \varepsilon \big( a_i + w_i(t) \big),
  \label{eq:initproto}
\end{equation}
with $a_i \sim \mathcal{N}(0,1)$ a per-rotor constant and $w_i(t)$ a
unit-variance Ornstein--Uhlenbeck process with 1\,s correlation time,
both frozen per seed; the single knob $\varepsilon \in \{0.1, 0.3, 0.5,
1.0, 2.0\}\revps$ sweeps initialization quality from near-oracle to
beyond every refiner's capture range. All refiners, ours and baselines,
receive identical draws. Blind operation is reported under two
protocols: (i) \emph{shared blind init}---every refiner starts from our
global stage output, isolating refinement quality; and (ii)
\emph{end-to-end blind}---each method uses its native front end (IAVKF:
dual-fast ridge detection per its paper; VKF+SWT: multiridge extraction;
tacholess order tracking: its own multi-order probabilistic seed).
Protocol (ii) is the headline.

\subsection{Baselines}
\textbf{IAVKF}~\cite{jiang2024iavkf}: faithful reimplementation (greedy
component peeling with energy-ratio stopping, $\alpha^0 = 10^8$, $\mu =
10^{-8}$ per the paper), made multi-rotor-fair in exactly one respect:
components are peeled per (rotor, harmonic) using our track table rather
than free spectral peaks; everything else is theirs.
\textbf{VKF+SWT}~\cite{li2020vkswt}: VK envelope extraction with external
synchrosqueezed-wavelet ridge IF. \textbf{Tacholess order
tracking}~\cite{peeters2019tacholess}: generalized-demodulation iterated
order spectra~\cite{olhede2005generalized}; our implementation is the
iterated-warp member of this family. \textbf{VK frequency loop}: the
inert classical loop of Sec.~\ref{sec:vkloop}, kept as an ablation row.
Neural pitch estimators are out of scope; rotor-conditioned
enhancement~\cite{gulli2025rotorcond} is the \emph{consumer} of this
output, not a competitor.

\subsection{Metrics}
PIT-MAE (\revps), pooled and per-rotor; the \emph{capture curve}, final
MAE versus $\varepsilon$; percent of windows converged (MAE $<
0.3\revps$); runtime per 16\,s window (single CPU core). On real data,
where telemetry is itself imperfect, the label-free reconstruction
residual $\mathcal{R}(r) = \|x - \hat x_r\|/\|x\|$
(Sec.~\ref{sec:vksolve}) serves as referee: it is monotone in trajectory
quality and consults no reference, so it can adjudicate differences
smaller than the reference's own noise.

\section{Results}
\label{sec:results}

\subsection{Synthetic Tiers}

Tables~\ref{tab:s1}--\ref{tab:s3} report the capture sweep of
Sec.~\ref{sec:initproto} on S1--S3. Measured anchors for our chain: on S1
the phase-increment pass holds an exact initialization to $0.055\revps$,
removes a $+1.0\revps$ offset to $0.107\revps$, and loses lock beyond
${\approx}1.5\revps$. On S2 the four-rotor 0\,dB floor is
${\approx}0.2\revps$, \emph{proved} by a sibling-oracle arm that
reconstructs and removes the other rotors from ground truth: oracle
$0.201$ versus blind $0.203$---additive noise dominates, and no sibling
scheme can do better. On S3 (measured $q_k$, $-5$\,dB) the tacholess warp
refiner reaches $0.04\revps$ and the phase-increment pass
$0.09$--$0.14\revps$ from a $+1.0\revps$ initialization.

\begin{table}[t]
  \centering
  \caption{S1 (single rotor, locked phases, $+10$\,dB): final PIT-MAE
    (\revps) vs.\ initialization error $\varepsilon$ (\revps), and blind.
    Ours measured; cells marked $\ast$ are \pending{S1 baseline run}.}
  \label{tab:s1}
  \scriptsize
  \setlength{\tabcolsep}{3pt}
  \newcommand{\pd}{{\color{gray}$\ast$}}
  \begin{tabular}{@{}lcccccc@{}}
    \toprule
    Method & $0.1$ & $0.3$ & $0.5$ & $1.0$ & $2.0$ & blind \\
    \midrule
    Ours (full) & 0.055 & \pd & \pd & 0.107 & $>1$ & \pd \\
    Ours w/o $\pi$-Kalman & \pd & \pd & \pd & \pd & \pd & \pd \\
    $\pi$-Kalman, shared init & 0.055 & \pd & \pd & 0.107 & $>1$ & --- \\
    IAVKF~\cite{jiang2024iavkf} & \pd & \pd & \pd & \pd & \pd & \pd \\
    VKF+SWT~\cite{li2020vkswt} & \pd & \pd & \pd & \pd & \pd & \pd \\
    Tacholess OT~\cite{peeters2019tacholess} & \pd & \pd & \pd & \pd &
      \pd & \pd \\
    VK loop (ablation) & \pd & \pd & \pd & \pd & \pd & \pd \\
    \bottomrule
  \end{tabular}
\end{table}

\begin{table}[t]
  \centering
  \caption{S2 (four rotors, twin pair $0.5\revps$, 0\,dB): layout as
    Table~\ref{tab:s1}; $\ast$ cells \pending{S2 baseline run}. The $0.2\revps$ floor is proved by the sibling
    oracle (0.201) meeting the blind chain (0.203).}
  \label{tab:s2}
  \scriptsize
  \setlength{\tabcolsep}{3pt}
  \newcommand{\pd}{{\color{gray}$\ast$}}
  \begin{tabular}{@{}lcccccc@{}}
    \toprule
    Method & $0.1$ & $0.3$ & $0.5$ & $1.0$ & $2.0$ & blind \\
    \midrule
    Ours (full) & \pd & \pd & \pd & \pd & \pd & 0.203 \\
    Sibling oracle & \multicolumn{5}{c}{---} & 0.201 \\
    IAVKF & \pd & \pd & \pd & \pd & \pd & \pd \\
    VKF+SWT & \pd & \pd & \pd & \pd & \pd & \pd \\
    Tacholess OT & \pd & \pd & \pd & \pd & \pd & \pd \\
    VK loop (ablation) & \pd & \pd & \pd & \pd & \pd & \pd \\
    \bottomrule
  \end{tabular}
\end{table}

\begin{table}[t]
  \centering
  \caption{S3 (four rotors, measured phase diffusion, $-5$\,dB): layout
    as Table~\ref{tab:s1}; $\ast$ cells \pending{S3 baseline run}. Measured anchors at $\varepsilon = 1.0$:
    tacholess warp 0.04, our $\pi$-Kalman 0.09--0.14.}
  \label{tab:s3}
  \scriptsize
  \setlength{\tabcolsep}{3pt}
  \newcommand{\pd}{{\color{gray}$\ast$}}
  \begin{tabular}{@{}lcccccc@{}}
    \toprule
    Method & $0.1$ & $0.3$ & $0.5$ & $1.0$ & $2.0$ & blind \\
    \midrule
    Ours (full) & \pd & \pd & \pd & 0.09--0.14 & \pd & \pd \\
    Tacholess OT (warp) & \pd & \pd & \pd & 0.04 & \pd & \pd \\
    IAVKF & \pd & \pd & \pd & \pd & \pd & \pd \\
    VKF+SWT & \pd & \pd & \pd & \pd & \pd & \pd \\
    VK loop (ablation) & \pd & \pd & \pd & \pd & \pd & \pd \\
    \bottomrule
  \end{tabular}
\end{table}

\subsection{Real Data}

Table~\ref{tab:real} reports the frozen 15-window protocol. Our measured
blind rows: on DREGON the chain ties the blind-VK baseline pooled ($1.825
\to 1.819\revps$)---the decoupling stage no longer subtracts precision,
and sits at or below baseline on 8 of 9 windows. On FLY124 cruise the
chain improves $2.418 \to 1.864\revps$ ($-23\%$) with cross-window
seeding, which pools a base across cruise windows of the same recording
and is off by default; without it, $2.380$. With a neural-predictor
initialization the refinement improves $1.016 \to 0.859\revps$ pooled,
the best non-oracle FLY124 figure.

\begin{table}[t]
  \centering
  \caption{Real data, pooled PIT-MAE (\revps), end-to-end blind. Ours
    measured; cells marked $\ast$ are \pending{real-data baseline runs}.}
  \label{tab:real}
  \scriptsize
  \setlength{\tabcolsep}{3pt}
  \newcommand{\pd}{{\color{gray}$\ast$}}
  \begin{tabular}{@{}lccc@{}}
    \toprule
    Method & DREGON steady & DREGON ramp & FLY124 cruise \\
    \midrule
    Blind-VK baseline & 1.825 & \pd & 2.418 \\
    Ours (gated chain) & 1.819 & \pd & 2.380 \\
    Ours ($+$ cross-window seed) & --- & --- & 1.864 \\
    Ours (neural init) & \pd & \pd & 0.859 \\
    IAVKF & \pd & \pd & \pd \\
    VKF+SWT & \pd & \pd & \pd \\
    Tacholess OT & \pd & \pd & \pd \\
    \bottomrule
  \end{tabular}
\end{table}

\begin{figure}[t]
  \centering
  \figplaceholder{6em}{F2: capture-radius arithmetic diagram}
  \caption{Capture-radius arithmetic. Fixed band: the radius $B/k$
    shrinks with harmonic order, collapsing on exactly the harmonics the
    weighting favors. $k$-scaled band $B_k = kB_0$: constant radius $B_0$
    at every order, while inter-rotor discrimination is preserved because
    the $k$-th teeth of two rotors are $k|r_i - r_j|$ apart.}
  \label{fig:capturearith}
\end{figure}

\begin{figure}[t]
  \centering
  \figplaceholder{6em}{F3: DREGON ramp spectrogram + tracked combs}
  \caption{DREGON takeoff ramp: whitened spectrogram with the four
    tracked combs overlaid.}
  \label{fig:ramp}
\end{figure}

\begin{figure}[t]
  \centering
  \figplaceholder{6em}{F4: capture curves, all methods}
  \caption{Capture curves: final PIT-MAE vs.\ initialization error
    $\varepsilon$ under the protocol of Sec.~\ref{sec:initproto}, all
    methods.}
  \label{fig:capturecurves}
\end{figure}

\begin{figure}[t]
  \centering
  \figplaceholder{6em}{F5: measured weight law (data exists:
    scripts/phase\_noise\_cov)}
  \caption{Measured per-harmonic information law. Inverse variance
    $1/v_k$ rises as $k^{2.0}$ (DREGON, 36 units) and $k^{1.5}$
    (Matrice~100, 52 units) at a fixed band, with no amplitude factor;
    under $B_k = 0.25k$ the exponent collapses to ${\approx}0$. The
    deployed weight $k^2|p_k|$ is effectively flat.}
  \label{fig:weightlaw}
\end{figure}

\begin{figure}[t]
  \centering
  \figplaceholder{6em}{F6: FLY124 twin-pair qualitative}
  \caption{FLY124 twin pair: pair-joint two-peak tracking through a
    collided band (qualitative).}
  \label{fig:twinpair}
\end{figure}

\subsection{One Pass, Not a Loop}

Table~\ref{tab:iterharm} measures repeated application of the
phase-increment pass on S1: error grows monotonically, because each pass
injects ${\approx}0.05\revps$ of fresh estimation noise that accumulates
in the track. This justifies the one-pass design of
Algorithm~\ref{alg:chain}.

\begin{table}[t]
  \centering
  \caption{Iteration harm: PIT-MAE (\revps) after $n$ repeats of the
    phase-increment pass (S1, measured).}
  \label{tab:iterharm}
  \footnotesize
  \begin{tabular}{@{}lccccc@{}}
    \toprule
    Repeats $n$ & 1 & 2 & 3 & 4 & 5 \\
    \midrule
    PIT-MAE & 0.167 & 0.196 & 0.226 & 0.244 & 0.268 \\
    \bottomrule
  \end{tabular}
\end{table}

\subsection{Does the Revived Loop Work?}
\label{sec:revival}

Section~\ref{sec:vkloop} proposed two fixes to the inert frequency loop:
the $k$-scaled band~\eqref{eq:kscaled} and the measured weight law.
Table~\ref{tab:revival} tests whether either revives it, iterating the
loop from controlled offsets.

\begin{table}[t]
  \centering
  \caption{Frequency-loop revival sweep: final PIT-MAE (\revps) from
    controlled offsets. Cells marked $\ast$ are \pending{job bash-07eb1a}.}
  \label{tab:revival}
  \footnotesize
  \setlength{\tabcolsep}{4pt}
  \newcommand{\pd}{{\color{gray}$\ast$}}
  \begin{tabular}{@{}lccc@{}}
    \toprule
    Arm & $0\revps$ & $0.3\revps$ & $1.0\revps$ \\
    \midrule
    Control (classical loop) & \pd & \pd & \pd \\
    $k$-scaled band & \pd & \pd & \pd \\
    IAVKF-adaptive band & \pd & \pd & \pd \\
    Both & \pd & \pd & \pd \\
    Both $+$ measured weights & \pd & \pd & \pd \\
    $\pi$-Kalman (reference) & \pd & \pd & \pd \\
    \bottomrule
  \end{tabular}
\end{table}

\subsection{Runtime}

The chain runs at ${\approx}0.9\times$ the runtime of its predecessor
pipeline on a single CPU core, dominated by blind seeding and the
mid-band DP stages; the banded envelope solves and the phase-increment
pass cost seconds per 16\,s window. Per-stage timings: \pending{exact
re-timing of the current chain, per stage, per 16\,s window}.

\section{Discussion and Limitations}
\label{sec:discussion}

\emph{What is established.} The capture-range analysis of
Sec.~\ref{sec:vkloop} is proved arithmetically and confirmed by three
independent instruments; the $0.2\revps$ four-rotor floor at 0\,dB is
proved by an oracle that the blind chain meets; the real-data gains are
measured on a frozen protocol. The baseline comparison and the loop
revival are designed but not yet run, so no claim of superiority over
IAVKF or tacholess order tracking is made beyond the structural argument.

\emph{The displaced comb.} On DREGON translating flight the surviving low
harmonics sit $0.3$--$0.5\revps$ \emph{below} $k$ times the telemetry
rate; a hover contrast is $3$--$4\times$ weaker, and a synthetic control
exonerates the estimator, implicating the sound field of translation. Any
audio-locked method is therefore charged $0.3$--$0.5\revps$ by a
telemetry-referenced metric on such data, and sub-$0.5\revps$ claims
should be read against that bias.

\emph{Assumptions.} The diffusion model takes $b_k$ Brownian with
Gaussian increments and time-constant $q_k$ within a pass; the
band-correlation factors in~\eqref{eq:incvar} assume ideal bands, with
corrections of order $10\%$. The S-tier ground truth is point-sampled
onto the frame grid without anti-alias filtering, which inflates
synthetic scores at high trajectory bandwidth---a benchmark artifact,
quantified and held fixed across methods.

\emph{Open.} The measured coherence-time budgets $\tau_k$ bound coherent
integration for any estimator; assembling them into a Cram\'er--Rao
bound for blind multi-rotor rate tracking is open.

\section{Conclusion}
\label{sec:conclusion}

Blind recovery of four rotor rates from quadrotor ego-noise is achievable
at twin separations below $1\revps$ by a staged pipeline whose stages are
matched by capture range: a global comb-DP stage that cannot be trapped,
an exact coupled VK envelope solve, and one phase-increment ML pass. The
classical VK frequency loop is inert at the blind operating point for a
stated arithmetic reason; published adaptive variants manage the same
radius implicitly, and the $k$-scaled bandwidth with measured
per-harmonic weights makes the schedule explicit. The initialization
protocol and the label-free referee make refiner comparisons
honest---including against the reference telemetry itself.

\bibliographystyle{IEEEtran}
\bibliography{references}

\end{document}
